\scanpage{95}
\mspage{94}
faut pour définir une extension
\[
1\mathbin{\text{---}}\pi\longrightarrow
\underbrace{\mathcal E_{\Phi_{\unsure{s}}}}_{\substack{\Phi\ \text{avec sa}\\(\pi,\mu)\text{ structure}}}
\longrightarrow\mu\mathbin{\text{---}}1
\]
\margin{\unsure{Avec les hyp. faites}, cette extension $\Rightarrow$ centrale,
mais on aurait pu faire opérer $\mu$ non trivialement sur $\unsure{\pi}$\ldots}
(\uline{NB} on n'a pas utilisé la commutativité de $\mu$
encore) --- et \inserted{inversement}, une \deleted{telle} extension
(\unsure{pas même|pas néc.} commutative)
\[
1\mathbin{\text{---}}\pi\mathbin{\text{---}}\mathcal E\mathbin{\text{---}}\mu\mathbin{\text{---}}1
\]
\margin{la commutativité de $\pi$ \unsure{entraîne une considérable simplification}}
plus la donnée d'un torseur $E$ sous $\mu$,
permet de construire canoniquement un
$\pi$-groupoïde $\Phi=\mathcal E(E)$, plus (par transport
de structure) une opération de $\mu$ sur
$\Phi=\mathcal E(E)$. En résumé, on a \deleted{une}
une équivalence de catégories
\diagram{\begin{tikzcd}[ampersand replacement=\&, column sep=large]
\begin{gathered}
\pi\text{-groupoïdes}\ \inserted{\text{conn. }\Phi},\ \text{avec}\\
\text{op. du groupe }\mu\text{ sur}\\
(\Phi,\unsure{E|R}),\ \text{tel que }\operatorname{Ob}\Phi\\
\text{soit un }\mu\text{-torseur}
\end{gathered}
\arrow[r,"\approx"] \&
\begin{gathered}\text{extensions de}\\\mu\text{ par }\pi\end{gathered}
\end{tikzcd}}\quad(17)

Ceci dit, l'extension (13) étant donnée
une fois pour toutes, on dit qu'on a
une $\mathcal E(\pi,\mu)$ structure sur $C'$ au dessus de
$C$, si \deleted{\unsure{pour tout}}
\inserted{(en plus des conditions a), b), c))}
pour tout $\unsure{\xi}\in\operatorname{Ob}C$ on se donne un iso. entre
